\section{思考与练习}

\begin{exercise}
切换到 buddy 后端（\texttt{KCONFIG\_PMM\_BACKEND=1}），运行 \texttt{pmm state}。
对比 bitmap 和 buddy 的 \texttt{free} 数量。buddy 为什么可能更少？
（提示：buddy 的元数据——\texttt{free\_list} 节点和 order 标记——也占物理页。）
\end{exercise}

\begin{exercise}
运行 \texttt{pmm alloc 100} → \texttt{pmm state} → \texttt{pmm freeall} → \texttt{pmm state}。
free 是否完全恢复？如果差了几页，用 \texttt{pmm dump} 查看哪些页没被释放。
\end{exercise}

\begin{exercise}
在 \texttt{pmm\_bitmap\_alloc\_page} 中注释掉 round-robin 逻辑
（\texttt{g\_region\_rr = ri}），改为总是从 region 0 开始扫描。
连续分配 1000 页后用 \texttt{pmm dump} 查看各 region 的分布。
和 round-robin 对比有什么不同？
\end{exercise}

\begin{exercise}
（挑战）bitmap 分配器不支持连续多页分配。实现一个
\texttt{pmm\_bitmap\_alloc\_pages(unsigned count)} 函数：
在位图中找到 \texttt{count} 个连续的 0 bit，全部置 1 并返回起始物理地址。
这比 buddy 的 order 分配更灵活，但扫描开销更大。
\end{exercise}

% ============================================================================

